%% ============================================================================
\section{Properties of the Score}
\label{app:props}
%% ============================================================================

Theorem~\ref{thm:mono} is the main result, and nothing in this section is needed to state or prove
it. The section collects what else is true of the score: which of its choices the guarantee depends
on (App.~\ref{app:why}), and the properties that make $\Lambda$ read as an amount of evidence
(App.~\ref{app:further}). Proofs are in App.~\ref{app:proof-why} and App.~\ref{app:proof-props}.

\subsection{Why This Score}
\label{app:why}

Theorem~\ref{thm:mono} rests on two choices, the relation tables and the world at which the MRF is
read, and neither can be changed.

\begin{theorem}[Sign-definiteness is necessary and sufficient]
\label{thm:necessity}
Fix the priors and $\xs$. Adding a positive table $\psi$ over a pair of claims $e$ never decreases
(resp.\ increases) $\Lambda$, for every relation graph, if and only if $\psi(\xs_e)=\max\psi$
(resp.\ $\min\psi$).
\end{theorem}

The tables of Eq.~\ref{eq:psi} take only the values $1$ and $1-p_r$, so they meet the condition at
every $\xs$ (App.~\ref{app:signdef}); a table whose committed cell lies strictly between its
extremes helps in some graphs and hurts in others. Required at every committed world, the condition puts every cell at the table's
maximum or minimum, so the table takes at most two values: a two-valued table.

\paragraph{Row-stochastic tables.}
\label{app:rowstoch}
A natural alternative to Eq.~\ref{eq:psi} uses tables that are
row-stochastic in the source, so that a false premise leaves its conclusion at the prior
(App.~\ref{app:earlier}). Such tables fail Theorem~\ref{thm:necessity}'s test. At $\xs=\one$ they
are sign-definite with the intended sign only for $p\ge\frac12$. Below it the committed cell of an
entailment is the table's minimum, so a weak entailment lowers the score and a weak contradiction
raises it. When the source is false at $\xs$, the committed cell holds the prior's value
$\frac12$, which is neither the maximum nor the minimum, so they fail at every weight except
$p=\frac12$, where the table is constant and has no effect. For weight changes,
App.~\ref{app:proof-why} gives the analogue of Theorem~\ref{thm:necessity} as a sufficient
condition.

\begin{proposition}[No other readout world]
\label{prop:unique}
For $n\ge2$ and any world $x\neq\xs$, some elementary improvement strictly lowers
$\log\PR(x)-\log\pi(x)$, for any priors and any $\R$.
\end{proposition}

The order and the readout must use the same world. Over \thmSteps{} random improving edits,
Shogenji's measure, which ignores $\xs$, falls on \fallsPctShogenjiSigndef\% and $\Lambda$ on none
(App.~\ref{app:verify}). So does the mean posterior marginal $\frac1n\sum_i\PR(a_i{=}1)$: deleting a
contradiction that holds down a hub claim, which itself contradicts five others, lowers it from
$\starFiveMMa$ to $\starFiveMMb$, and it falls on \fallsPctMmSigndef\% of the random edits.

\paragraph{The readout world.} Proposition~\ref{prop:unique} does not say that $\xs$ is the right
world; choosing it is the factuality pipeline's job. It says that the order and the readout must
use the same world. It also does not exclude readouts that combine several worlds; of the five
App.~\ref{app:earlier} examines, four are not monotone and the fifth violates A2.

\paragraph{Why not count relations?} Monotonicity alone admits trivial scores: the signed weight sum
$\sum_{\Rsup}p_r-\sum_{\Rcon}p_r$ is monotone by construction. But a count credits every relation in
full, whether or not it rules out anything plausible, so it depends on how an argument is presented
rather than on what it rules out. Two entailments $a\Rightarrow b$ and $b\Rightarrow a$ of weight
$p$ multiply to exactly the table of one equivalence of weight $p$, so they induce the same $\PR$
and the same $\Lambda$, but they count $2p$ against $p$, which violates A2. Repeating a support
likewise adds its weight each time, where $\Lambda$ saturates (Prop.~\ref{prop:props}(vii)).
Averaging removes the dependence on the number of relations but breaks monotonicity, since a weak
support lowers an average. $\Lambda$, by contrast, weighs each relation by what it rules out under
the rest of the model. Proposition~\ref{prop:decomp} locates the sum: the first-order expansion of $\Lambda$ is the
linear evidence $\Llin$ (Eq.~\ref{eq:llin}), and the sum is $2\Llin$ with every relation taken to
forbid half the prior mass of its pair, which removes the asymmetry between supports and conflicts.
Both drop the interaction term $I_\R$, where the rest of A2 lives. Normalizing the sum to $[-1,1]$ would bound it but break A1, since a further satisfied
support leaves a fully satisfied graph unchanged; $\LCS$ is bounded and strictly monotone at once.

\paragraph{What is not forced.} These results fix the class of tables and the readout world; they
do not make $\Lambda$ the unique score satisfying A1--A3. Any score that is strictly increasing in
$\Lambda$ for a fixed prior and takes a prior-independent value at $\Lambda=0$ satisfies them too,
which is how $\LCS$ does. And any reweighting $p\mapsto1-(1-p)^\beta$ with $\beta>0$ gives another
monotone score of the same form, where different $\beta$ can rank incomparable graphs differently.
We take $\beta=1$ because it makes $p_r$ the probability that a check catches a violation, so the
check reading holds with the weights as estimated.

\subsection{Further Properties}
\label{app:further}

\begin{proposition}[Further properties]
\label{prop:props}
\label{prop:basic}
\label{prop:neg}
\label{prop:sat}
(i) $\Lambda(\varnothing)=0$ for every $\pi$. (ii) $\Lambda\ge0$ if $\R$ has only supports and
$\Lambda\le0$ if only conflicts. (iii) If $\R_1$ and $\R_2$ have scopes in disjoint sets of claims,
$\Lambda(\R_1\uplus\R_2)=\Lambda(\R_1)+\Lambda(\R_2)$. (iv) Replacing a claim by its negation,
together with its verdict, its prior and every coupling at it, preserves every sign, $\Lambda$ and
$\LCS$. (v) $\sum_{r\in\Rcon}\log(1-p_r)\le\Lambda(\R)\le-\log\pi(\xs)$. (vi) A support's gain is
at most $-\log\bigl(1-p_r(1-\PR(\xs))\bigr)$. (vii) For $p_r>0$, a second copy of a support adds
strictly less than the first.
\end{proposition}

$\Lambda$ is thus evidence, in nats, that adds across independent sub-arguments, and redundant
support saturates. Item by item:
\begin{itemize}[leftmargin=1.2em,itemsep=1pt,topsep=2pt]
\item \emph{Neutrality} (i). With no relations $\PR=\pi$, so $\Lambda=0$ and $\LCS=\frac12$ whatever
  the priors; this is A3.
\item \emph{Pure graphs} (ii). An argument that draws only relations its committed world upholds
  cannot score below neutral, and one that draws only relations it breaks cannot score above it.
\item \emph{Additivity} (iii). Sub-arguments over disjoint claims contribute independently, so
  $\Lambda$ is an amount of evidence in \emph{nats} (natural-log units: a change of $\Delta$ nats
  multiplies $\PR(\xs)$ by $e^{\Delta}$) that adds across independent parts of an argument. $\LCS$ is
  not additive (App.~\ref{app:lcs-lambda}).
\item \emph{Negation invariance} (iv). Whether a decomposition records a denied claim as itself or
  as its asserted negation does not matter, provided the verdict, the prior and the couplings are
  rewritten consistently (App.~\ref{app:framework}, \emph{Claims and their negations}).
\item \emph{Bounds} (v). Since every factor is at most $1$, $Z_\R\le1$ and $-\log Z_\R\ge0$, so the
  conflicts alone set a floor: $\Lambda$ cannot fall below the sum of their $\log(1-p_r)$, whatever
  the supports do. The ceiling comes from $\PR(\xs)\le1$; with weights below $1$ it is approached
  only as the relations come close to forcing $\xs$. The floor tends to $-\infty$ as a conflict's
  weight tends to $1$ (App.~\ref{app:weights}).
\item \emph{Bounded gain} (vi). A support earns credit only for the probability it removes from
  worlds other than $\xs$: in Lemma~\ref{lem:edit}, $q_r\le1-\PR(\xs)$. Once the relations nearly
  force $\xs$, further support earns almost nothing.
\item \emph{Saturation} (vii). A repeated support earns less each time, and repeated copies converge
  rather than accumulate (App.~\ref{app:proof-props}), unlike a count. A repeated \emph{conflict} does the opposite: each further copy costs more than the first (at
$p=\dcP$, a violated \con{} costs $\dcCostOne$ nats alone), its marginal cost rising to
$-\log(1-p)$, so the total cost grows without bound, because $I_\R>0$ lowers $\Lambda$ whatever the
signs (\S\ref{sec:decomp}).
\end{itemize}
Each relation's leave-one-out contribution $\Delta_r=\Lambda(\R)-\Lambda(\R\setminus\{r\})$, the
bars of Figure~\ref{fig:voyager}, is Lemma~\ref{lem:edit} applied to $\R\setminus\{r\}$, so it has a
guaranteed sign: positive for a support and negative for a conflict, whenever $p_r>0$. The contributions are not
additive and need not sum to $\Lambda$, since each is computed against the rest of the model.
Weighted mini-bucket elimination \citep{wmb}, for graphs beyond exact inference, preserves the
guarantees only up to the approximation.
